\chapter{Methodology}

This section presents a five-step procedure for implementing cdt:ucum in an RDF library. The procedure is general. It applies to any programming language, any RDF library that supports custom datatypes, and any UCUM backend available in that language. It assumes no particular library architecture and no particular level of SPARQL support. Both vary across environments, and the procedure is designed to make that variation explicit rather than to hide it.

Each step states its goal, prescribes the actions the implementer must carry out, and names the output it produces. The output of each step is the input of the next. This chaining is what makes the procedure repeatable: a second implementer starting from the same specification and the same candidate libraries should reach the same decisions for the same reasons.


\begin{figure}[ht]
\centering
\begin{tikzpicture}[
    stepbox/.style={
        rectangle, draw, thick,
        rounded corners=4pt,
        minimum width=9.0cm,
        minimum height=1.0cm,
        text width=8.4cm,
        align=left,
        inner sep=8pt,
        font=\small
    },
    decnode/.style={
        diamond, draw, thick,
        minimum width=2.8cm,
        minimum height=0.9cm,
        align=center,
        font=\small,
        aspect=2.8
    },
    haltnode/.style={
        rectangle, draw, thick,
        rounded corners=14pt,
        fill=gray!20,
        minimum width=3.6cm,
        minimum height=0.9cm,
        align=center,
        inner sep=6pt,
        font=\small
    },
    arr/.style={-{Stealth[length=2.5mm,width=2mm]}, thick},
    lbl/.style={font=\footnotesize},
]

\node[stepbox] (s1) {%
    \textbf{Step 1:} Identify Candidate Libraries\\[4pt]%
    {\footnotesize RDF Library\enspace|\enspace UCUM Backend}%
};

\node[haltnode, right=3.0cm of s1] (stop) {Language not viable\\(stop)};

\node[stepbox, below=0.9cm of s1] (s2) {\textbf{Step 2:} Analyse Extension Points and APIs};
\node[stepbox, below=0.7cm of s2]  (s3) {\textbf{Step 3:} Assess Feasibility};
\node[decnode,  below=0.7cm of s3]  (gng) {Go\,/\,No-go};
\node[stepbox, below=0.7cm of gng] (s4) {\textbf{Step 4:} Implementation};
\node[stepbox, below=0.7cm of s4]  (s5) {\textbf{Step 5:} Validation};

% main flow
\draw[arr] (s1) -- (s2);
\draw[arr] (s2) -- (s3);
\draw[arr] (s3) -- (gng);
\draw[arr] (gng.south) -- node[lbl, right=1.5mm] {go} (s4.north);
\draw[arr] (s4) -- (s5);

% exit branch
\draw[arr] (s1.east) -- node[lbl, below=1.5mm] {either unavailable} (stop.west);

% no-go loop back
\draw[arr] (gng.west) -- node[lbl, below=1.5mm] {no-go} ++(-3.2,0) |- (s1.west);

\end{tikzpicture}
\caption{Five-step procedure with exit at Step~1 and go/no-go branch after Assess Feasibility}
\label{fig:five-step-procedure}
\end{figure}

\section{Identify Candidate Libraries}

The goal of this step is to identify two separate things: an RDF library to host the implementation, and a UCUM backend to handle unit parsing and arithmetic. The two searches are independent and proceed in parallel.

For the RDF library, the implementer must list the libraries available in the target language and subject each to a surface check covering three properties. The first is active maintenance, which is the floor rather than the whole of it: the library's age, contributor base, and breadth of adoption count alongside it, and a library may lack a major feature and still be the natural host if the community has settled on it. The second is SPARQL coverage; a library without a SPARQL engine is not discarded outright, but its available query mechanisms must be documented, because the operator obligations may still be deliverable in reframed form. The third is the presence of a custom datatype mechanism; without it, the registration obligation cannot be met.

For the UCUM backend, the implementer must identify libraries that can parse UCUM codes, convert between units, and perform arithmetic. Since UCUM-native libraries are scarce in most languages, the search must also include general unit-arithmetic libraries; these cannot serve alone but can pair with a UCUM-native parser. When a combination is chosen, the implementer must record which library handles which responsibility, because that boundary shapes the design of the internal quantity representation; two of the four implementations required this arrangement.

The step can produce a null result on either side. If no RDF library with a custom datatype mechanism exists for the target language, or no library or pairing can parse UCUM expressions and perform arithmetic at runtime, the shortlist on that side is empty. An empty shortlist on either side means implementation is not possible and the procedure stops.

The output is a shortlist of each kind, ordered by suitability, with the reason for every elimination recorded. The ordering matters because a no-go at the feasibility step returns to the next candidate on the shortlist rather than restarting the search.

\section{Analyze Extension Points and APIs}

The goal of this step is to read the internals of both shortlisted libraries and establish how each obligation of the specification maps onto what they actually expose.

For the RDF library, the analysis covers three areas. The first is datatype registration: how the library associates a datatype IRI with a custom type, what happens at parse time when a literal carrying that IRI is encountered, whether the association can be made from outside the core, and how the extension becomes active. The second is SPARQL operator evaluation: whether each evaluation point for equality, comparison, and arithmetic is exposed for extension or hardcoded, and when hardcoded, whether the host language permits runtime intervention from outside the library; the answer determines whether a workaround is available or the obligation requires a fork. The third area is the path from SPARQL constructs to those evaluation points: whether FILTER and BIND share the same machinery as ORDER BY and the aggregates SUM and AVG, or reach separate handlers, since the answers determine whether one hook covers everything or several are needed. Finally, the implementer must check whether the library provides a function registry invocable by IRI from inside a query; the sameDimension obligation depends on it.

For the UCUM backend, the analysis covers four areas. The first is the parsing API: what it accepts and returns. The second is arithmetic: whether the backend converts between units, tracks dimensions through operations, and signals incommensurable quantities. The third is the numeric type; a floating-point backend sets a precision ceiling no outer layer can raise. The fourth is coverage: whether all UCUM codes are supported or gaps exist.

The two cases examined here show why the analysis cannot stop at the library level. In Java, adjacent operators in the same engine received opposite answers; in Python, the operator module was reachable but ordering and aggregation held their own captured references, making an intervention appear complete while leaving those paths untouched.

The output of this step is a map from each obligation to the mechanism, if any, through which the library can satisfy it, together with the backend's capabilities and known gaps. This map is the sole input of the feasibility assessment that follows.

\section{Assess Feasibility}

The goal of this step is to decide, before any code is written, whether the candidate pair can satisfy the specification and to record how. The assessment runs through the obligations of Section 3.4 one by one.
One constraint frames the whole assessment: the implementation must be delivered as a standalone extension, ranking above every individual obligation. A fork produces a modified distribution users must run in place of the library they depend on and carry forward against every upstream release; Section 1 identified this as the shortcoming of the original Jena-ucum implementation. An obligation requiring a fork is recorded as not achievable.


For each obligation the assessment must reach one of three outcomes: satisfiable through a public API, satisfiable through a workaround, or not achievable without a fork. To reach the outcome, the implementer asks whether an extension point exists; if so, whether it is a documented public API or internals the library never promised to keep; and if not, whether the host language offers an alternative mechanism. The distinction matters because the cases age differently: a public API obligation survives upgrades, while a workaround must be re-verified at every upgrade and its interception points documented at the moment of choice.

The Python implementation shows the middle case. Comparison and arithmetic were not reachable through any public API: the engine validates operands against a fixed set of numeric types. The host language, however, permits replacing function references at runtime, so both were recorded as workarounds requiring re-verification at each library release.

When an obligation is not achievable, the implementer must either reframe it, delivering the same capability through a different interface, or return to Step 1 for the next candidate from the shortlist. The choice rests on three considerations: the centrality of the lost obligation, the availability of a viable shortlist alternative, and the library's weight in its ecosystem. A reframing must document what it achieves and gives up against the original wording; it is a documented deviation, not a fulfilment, and that record is what keeps the conformance claim honest.

The JavaScript implementation shows how ecosystem weight settles the choice. The library carries only a minimal query layer, placing the SPARQL operator obligations out of reach, yet it remains the library its community uses. 

The output of this step is a feasibility record stating, for each obligation, the case it falls under, the mechanism to be used, and any reframing adopted, together with the overall decision.


\section{Implementation}

The goal of this step is to turn the feasibility record into a working extension. The work follows a fixed order because each part depends on the one before it: the quantity representation first, registration second, the SPARQL hooks third, and the custom function fourth.

The implementer must begin with a class that wraps the UCUM backend and exposes every operation the value space requires: parsing a lexical form and writing one back, equality, comparison, arithmetic, unit conversion, and the dimension test. All backend access must pass through this class, keeping the rest of the extension independent of its API. The class must realise the value space defined in Section 3.4: two quantities in different units of the same dimension that denote the same value must compare as equal. Where the first step selected a pairing rather than a single backend, the division of labour recorded there becomes the internal structure of this class; the parser produces what the arithmetic library consumes.
With the representation in place, the implementer must register cdt:ucum and cdt:ucumunit through the extension point recorded in the feasibility record, after which the library parses a CDT literal into a typed object rather than an opaque string. 

The SPARQL hooks come next. The implementer must intercept equality, comparison, and arithmetic through the mechanisms recorded per obligation, and every handler must follow the same pattern: check whether the operands are CDT literals, evaluate them through the representation if they are, and pass the call to the original behaviour unchanged if they are not. The fallback matters as much as the interception, because the extension must leave the engine's treatment of every other datatype exactly as it found it. Ill-typed literals must be caught before any value operation and routed into the engine's error path, so that the affected solution is excluded, as Section 2.3 described, rather than the whole query failing. The map from the analysis step states how many interception points this takes and where none is needed. In two implementations, registration alone already delivered SPARQL equality. In the Python case, the interception reached twenty-five points across four modules.

The custom function completes the operator work. The implementer must register cdt:sameDimension through the function registry recorded in the analysis. The function extracts two quantities, asks the backend whether their dimensions match, and returns a boolean.

Implementation surfaces five cases that the analysis cannot show in advance and that must be expected in any new implementation. The first is units the backend cannot parse, with the special units of Section 3.3 first among them. The second is comparison across dimensions inside ORDER BY; the query must never abort on such data, and since SPARQL leaves this ordering undefined, any deterministic grouping or fallback is conformant. The third is unit preservation in SUM, AVG, MIN and MAX, where the engine's own accumulators work on bare numbers and would return a result stripped of its unit. The fourth is precision at the numeric extremes, where the backend's internal number type sets a ceiling no outer layer can lift. The fifth is the dimensionless case: a bare number must be read as a quantity with the unit one. Each case must be handled explicitly and the chosen behaviour written down, because several will return in the next step as expected failures.

The output of this step is the working extension together with the record of edge cases and the behaviour chosen for each. That record joins the feasibility record as input to the validation step.

\section{Validation}

The goal of this step is to verify that the extension satisfies each obligation it claims and to document what it does not. The output is not a count of passing tests but an account of conformance in which every claim is backed by a test and every limitation is visible.

The implementer must first test each obligation of Section 3.4. For lexical parsing, both valid and invalid inputs must be tested: a valid UCUM expression must parse to the correct value, and an invalid one must be rejected as ill-typed. Value equality is tested across units of the same dimension, including derived units and the dimensionless case. Comparison is tested within one dimension and across dimensions, where the error required by the specification must be raised. Arithmetic is tested over its full type coverage, including scalar operands on both sides of the operator. The dimension test is checked with compatible and incompatible pairs.
Each obligation must be tested at the level at which the feasibility record claims it. An obligation claimed at the query level requires tests that submit a query and inspect the results: a correct method the engine never calls passes every direct test and still fails the claim. Where an obligation was reframed, the tests target the reframed interface; no test may suggest the original form is available. The Python implementation tests at both levels: the quantity class directly, and the operators, ordering, and aggregates through queries. The Java implementation exercises equality through its interface but submits no query, so the query-level behaviour rests on the engine's documented routing rather than a test of this extension. The difference between these two positions is exactly what this step must make visible.

The implementer must then test the edge cases recorded during implementation. A query over data containing an ill-typed literal must complete, with the affected solutions excluded. Ordering over mixed dimensions must complete and follow the behaviour chosen for it. The aggregates must preserve units where the implementation claims they do. Precision is tested at scales far above and below the range of ordinary values, where the backend's numeric ceiling either holds or breaks. A literal written out and parsed back must keep its value. Activating the extension more than once must change nothing.

What the feasibility record marked as not achievable must now appear as expected failures. For each limitation, the implementer must write the test it causes to fail, mark it as expected, and attach the reason stating whether the cause is an inherited dependency limit or a gap in the extension itself; the two call for different responses from anyone who later builds on the work. A limitation carried this way is checked on every run and announces itself when a dependency upgrade removes it; one left out of the suite stays silent. The same two families recur across all four implementations: the special units of Section 3.3, whose conversion falls outside the algebra every backend implements, and precision loss where a backend computes in a bounded numeric type. A new implementation should expect to document both.

The output is the test suite together with the conformance account it expresses: the obligations met, the level at which each is met, the reframings, and the expected failures with their reasons. That account makes the extension's claims checkable by someone who did not build it, and is the form in which the four implementations are reported in Section 6.